\scanpage{85}
% archive foliation: 84; reused printed/typescript excerpt, oppositely oriented

où $p_{X'}$ désigne le morphisme structural $X'\quad X$.
D'autre part, pour donner un sens à 4.1. (ii) et justifier sa démonstration
indiquée de 4.1., il y a lieu d'introduire un univers \quad tel que \quad
\unsure{,}
et considérer les faisceaux intervenant dans les formules \unsure{4.1.}
\unsure{(ii)}
comme des faisceaux à valeurs dans \quad -(Ens).
% Blank spaces in the copy at the structural morphism and universe notation;
% no arrow or universe letter has been supplied from mathematical expectation.

\uline{Proposition \unsure{4I1.|4.1.}} (i) Le foncteur
$u^{\unsure{X|*}}$ est pleinement fidèle, donc pour tout faisceau $G$
sur $X_{\text{ét}}$, l'homomorphisme canonique
\[
G\qquad u_{\unsure{X|*}}u^{\unsure{X|*}}(G)
\]
est un isomorphisme.

(ii) Soient $G$ (resp. $F$) un faisceau abélien sur $X_{\text{ét}}$
(resp. $(\mathrm{Sch})/X$).
Alors les homomorphismes canoniques (définis par exemple comme
edge-homomorphismes de la suite \unsure{spectrahe|spectrale} de Leray V 5.5.1.)
suivants \unsure{sontx|sont} des isomorphismes:
\[
H^{\unsure{*}}(X_{\text{ét}},F\,u)\simeq
H^{\unsure{*}}((\mathrm{Sch})/X,F)
\]
\[
H^*(X_{\text{ét}},G)\simeq
H^{\unsure{*}}((\mathrm{Sch})/X,u^{\unsure{*}}(G)).
\]
En effet, le foncteur $u$ est un comorphisme (VI 7.3.6.) et on peut alors
appliquer VI 7.3.7. et VI 7.3.9.
